\section{Recomendaciones}

Con base en los resultados transitorios y anal\'iticos, se consolidan las siguientes recomendaciones de car\'acter normativo y metodol\'ogico (ver Tabla~\ref{tab:recomendaciones_resumen}), dirigidas a la gesti\'on de operaciones, el mantenimiento de integridad y las pr\'acticas de despacho log\'istico del tanque Globe~1130.

\begin{table}[H]
	\centering
	\caption{Consolidado de directrices y recomendaciones estrat\'egicas.}
	\label{tab:recomendaciones_resumen}
	\setcounter{itemcount}{0}
	\begin{tabularx}{\textwidth}{@{}lX@{}}
		\toprule
		\textbf{\'Area de Gesti\'on} & \textbf{Acci\'on Recomendada y Justificaci\'on} \\
		\midrule
		\multirow{3}{*}{\textbf{Termodin\'amica}}
		& \textbf{Incremento de Set-point del Flujo a 75°C}: Constituye la resoluci\'on econ\'omica m\'as viable, logrando contracci\'on del 52\% frente a tiempos requeridos sin modificar bomba o agitar el lecho. \\
		& \textbf{Instalaci\'on del Aislamiento T\'ermico Recomendado}: Instalar lana mineral de 50.8~mm (2'') sobre la chaqueta de media ca\~na y su tuber\'ia de servicio, seg\'un el espesor calculado en este estudio. Una vez instalado, programar inspecciones peri\'odicas mediante termograf\'ia infrarroja para detectar p\'erdidas localizadas por humedecimiento, da\~no mec\'anico o desprendimiento del recubrimiento exterior. \\
		& \textbf{Distribuci\'on Multi-ramal en Espiral (Flujo Paralelo)}: Reconfigurar estructuralmente la inyecci\'on continua del circuito hidrot\'ermico transmutando de recorrido \'unico a multi-entradas y multi-salidas. Este redise\~no mitiga la vasta ca\'ida de temperatura al final de la ruta por p\'erdida de sensiblilidad espec\'ifica y empodera el suministro garantizando un gradiente t\'ermico $\Delta T$ equitativo alrededor de la pared completa del cilindro. \\
		\midrule
		\multirow{3}{*}{\textbf{Log\'istica}}
		& \textbf{Pauta Anticipada de Lotes (192 Toneladas)}: Precalentar resguardando al menos 5.0h de anticipaci\'on pre-entrega log\'istica para suprimir riesgos de solidificaci\'on o estresamiento por viscosidad m\'axima de producto interno. \\
		& \textbf{Configuraci\'on de Descargas Consecutivas}: Agendar recepciones a tazas l\u00edneas contiguas de 1.5 horas perimetrales; evada los tiempos muertos considerando el incremento end\'ogeno continuo de $\sim 0.9$°C a partir de las 57°C org\'anicas. \\
		& \textbf{Fiscalizaci\'on Visual Fina}: Limitar y certificar que descargas tard\'ias continuas no transgredan o vulneren m\'aximas tolerables (como los detectados $> 68$°C al fin de la ruta multipunto). \\
		\midrule
		\multirow{3}{*}{\textbf{Integridad API}} 
		& \textbf{Bloqueo Activo a Techo Hidrost\'atico en $>$90\%}: Instituir bloqueo mec\'anico y sonoro en cotas l\'imite del 90\%; el lecho toriesf\'erico falla est\'aticamente por compresi\'on en factor marginal superior al CA previsto, proyectando riesgo real. \\
		& \textbf{Programa de Inspección Acelerado (M\'ax. 2 a\~nos)}: Debido al riesgo de \textbf{Corrosión-Fatiga} en el nuevo fondo (6.0 años de vida útil proyectada), se debe implementar un monitoreo mediante Ultrasonido de Arreglo de Fases (PAUT) y Mapeo de Corrosión con una frecuencia no mayor a 2 años. El enfoque principal debe ser la detección temprana de grietas por fatiga en la zona del nudillo y puntos de anclaje. \\
		& \textbf{Protocolos Sismo-receptivos}: Emisi\'on de an\'alisis base ante siniestros con $M>5.0$, priorizando el anillo as\'entrico dado el diferencial inmenso de carga gravitatoria (989 kN). \\
		\midrule
		\textbf{Integridad FEA}
		& \textbf{Refuerzo Estructural Proactivo en Zonas Críticas}: Basado en el mapeo de esfuerzos por FEA (ANSYS), se requiere la intervención estructural de los sectores donde el Factor de Seguridad (FS) sea inferior a 1.5. Es imperativo priorizar el nudillo fondo-cilindro y los anclajes de la chaqueta, donde se detectaron valores críticos de \textbf{FS $\approx$ 0.71} al 100\% de llenado. Se recomienda la implementación de anillos de rigidización perimetrales o placas de refuerzo (pads) con el fin de restaurar los márgenes normativos de ASME VIII Div. 1 y posibilitar la operación a plena capacidad operativa (100\% de llenado). El diseño detallado de este refuerzo debe desarrollarse en una fase posterior mediante un modelo FEA actualizado que valide la efectividad de las geometrías de refuerzo propuestas. \\
		\bottomrule
	\end{tabularx}
\end{table}

Las metodolog\'ias enunciadas en la tabla aseguran las cotas normativas y protegen activamente la prolongaci\'on de las 4 d\'ecadas log\'isticas esperadas para el contenedor en condici\'on actual.

\newpage
